\section*{第 \ref{ch:b1:spherical-astronomy} 章——球面天文学}

\begin{solution}{exo:b1:spherical-astronomy:1}
\omterm{def:b1:spherical-astronomy:celestial-sphere}{天极}的高度等于\omterm{def:b1:spherical-astronomy:geographic}{纬度}：\ang{52}。若
$\delta > \ang{90} - \ang{52} = \ang{38}$，则恒星永不落下，且若
$\delta < \ang{-38}$，则永不升起。
\end{solution}

\begin{solution}{exo:b1:spherical-astronomy:2}
$\qty{6.75}{h}\times\ang{15} = \ang{101.25}$. $H = \mathrm{LST} - \alpha =
\qty{10}{h} - \qty{6}{h}\,\qty{45}{min} = \qty{3}{h}\,\qty{15}{min}$，即
\ang{48.75}。\omterm{def:b1:spherical-astronomy:equatorial}{时角}为正：恒星在子午线以西，
它于三又四分之一个恒星小时前过子午线。
\end{solution}

\begin{solution}{exo:b1:spherical-astronomy:3}
$h = \ang{90} - \varphi + \delta_\odot$：夏至日为 \ang{73.44}，
春分、秋分时为 \ang{50}，冬至日为 \ang{26.56}。
\end{solution}

\begin{solution}{exo:b1:spherical-astronomy:4}
1900 年：能被 100 整除，但不能被 400 整除，不是闰年。2000 年：能被
400 整除，是闰年。2100 年：不是闰年。儒略年比\omterm{def:b1:spherical-astronomy:tropical-year}{回归年}
多出 $365.25 - 365.2422 = 0.0078$ 天，每世纪 \num{0.78} 天；格里高利年
则多出 $0.0003$ 天，每世纪 \num{0.03} 天。
\end{solution}

\begin{solution}{exo:b1:spherical-astronomy:5}
$\cos\Delta = \sin^2\ang{45} + \cos^2\ang{45}\cos\ang{90} = 0.5$，于是
$\Delta = \ang{60}$，距离为 $6371\times\pi/3 = \qty{6672}{km}$。
沿纬线，即半径为 $R\cos\ang{45}$ 的\omterm{def:b1:spherical-astronomy:spherical-triangle}{小圆}，距离为
$6371\times\cos\ang{45}\times\pi/2 = \qty{7076}{km}$。\omterm{def:b1:spherical-astronomy:spherical-triangle}{大圆}
的距离短约 \qty{404}{km}。
\end{solution}

\begin{solution}{exo:b1:spherical-astronomy:6}
$\sin h = \sin\ang{48}\sin\ang{20} + \cos\ang{48}\cos\ang{20}\cos\ang{30}
= 0.2542 + 0.5446 = 0.7987$, $h = \ang{53.0}$.
$\cos\mathrm{A} = (\sin\ang{20} - \sin\ang{48}\sin\ang{53.0})/(\cos\ang{48}
\cos\ang{53.0}) = (0.3420 - 0.5935)/0.4027 = -0.6246$, $\arccos = \ang{128.7}$;
恒星位于子午线以西（$H > 0$），故 $\mathrm{A} = \ang{360} -
\ang{128.7} = \ang{231.3}$：朝西南方向，高出地平线 \ang{53}。
\end{solution}

\begin{solution}{exo:b1:spherical-astronomy:7}
$\cos H_0 = -\tan\ang{60}\tan\ang{23.44} = -1.7321\times0.4335 = -0.7510$,
$H_0 = \ang{138.67}$，日长 $2\times138.67/15 = \qty{18.49}{h}$，约
\qty{18}{h}~\qty{29}{min}，与北纬 \ang{60} 曲线在夏至日的读数相符。
\end{solution}

\begin{solution}{exo:b1:spherical-astronomy:8}
中天时刻每天提前 \qty{3.93}{min}，因此 $30\times3.93 =
\qty{118}{min}$，即 30 天后，恒星约于 20:02 过子午线。要在
正午过子午线，须提前 \qty{10}{h} = \qty{600}{min}，所需时间为
$600/3.93 \approx 153$ 天。
\end{solution}

\begin{solution}{exo:b1:spherical-astronomy:9}
(a) 春分时正午高度角为 \ang{45}：$I = 1361\times\sin\ang{45}
= \qty{962}{W.m^{-2}}$。(b) 夏至日 $h = \ang{68.44}$，且
$(a/r)^2 = (149.598/152.100)^2 = 0.9674$:
$I = 1361\times0.9674\times\sin\ang{68.44} = \qty{1224}{W.m^{-2}}$。
较高的太阳足以抵消距离增大的影响。
\end{solution}

\begin{solution}{exo:b1:spherical-astronomy:10}
$\sin\delta_\odot = \sin\ang{23.44}\sin\ang{60} = 0.3445$,
$\delta_\odot = \ang{20.15}$. $\tan\alpha_\odot = \cos\ang{23.44}\tan\ang{60}
= 1.5890$，$\alpha_\odot = \ang{57.82}$（与 $\lambda_\odot$ 同象限）。
\omterm{def:b1:spherical-astronomy:equatorial}{赤经}落后 \ang{2.18}，即 $2.18\times4 \approx
\qty{8.7}{min}$ 的时间。沿\omterm{def:b1:spherical-astronomy:ecliptic}{黄道}匀速运动的太阳，
其在\omterm{def:b1:spherical-astronomy:equatorial}{赤经}上的运动仍不均匀，而\omterm{def:b1:spherical-astronomy:equatorial}{赤经}决定\omterm{def:b1:spherical-astronomy:equatorial}{时角}，
进而决定日晷：春分到下一个至日之间落后于平太阳，
至日到下一个春分之间超前。这正是
\omterm{def:b1:spherical-astronomy:solar-time}{时差}中由\omterm{def:b1:spherical-astronomy:ecliptic}{黄赤交角}引起的一项。
\end{solution}

\begin{solution}{exo:b1:spherical-astronomy:11}
升起时 $h = 0$，由\cref{prop:b1:spherical-astronomy:transform}可得
$\sin\delta = \cos\varphi\cos\mathrm{A}_0$；夏至日
$\delta = \varepsilon$，于是 $\cos\mathrm{A}_0 = \sin\varepsilon/\cos\varphi$。
当 $\mathrm{A}_0 = \ang{45}$ 时：$\cos\varphi = \sin\ang{23.44}/\cos\ang{45} =
0.5626$，北纬 $\varphi = \ang{55.8}$。当\omterm{def:b1:spherical-astronomy:geographic}{纬度}高于 $\ang{90} - \varepsilon =
\ang{66.56}$ 时，$\sin\varepsilon/\cos\varphi > 1$，不存在
升起\omterm{def:b1:spherical-astronomy:horizontal}{方位角}，因为当天太阳不落，即出现极昼。
\end{solution}

\begin{solution}{exo:b1:spherical-astronomy:12}
(a) 夹角为 $\theta$ 的球面二角形占球面的比例为 $\theta/2\pi$：
$\frac{\theta}{2\pi}\times4\pi R^2 = 2\theta R^2$。(b) 三对
球面二角形的总面积为 $2\times(2A + 2B + 2C)R^2$；它们恰好覆盖球面一次，
但三角形及其对径三角形各被覆盖三次，即
比所需多两次：$4(A + B + C)R^2 = 4\pi R^2 + 4S$，其中 $S$ 为
三角形的面积。由此 $S = R^2(A + B + C - \pi)$。(c) 三个直角：
$S = R^2(3\pi/2 - \pi) = \pi R^2/2$，为球面的八分之一（即一个卦限）。
图中三角形：$89.9 + 64.9 + 67.4 - 180 = \ang{42.2}$，即
\num{0.736} 弧度：$S = 0.736\,R^2$，占球面的 \qty{5.9}{\percent}。
\end{solution}

\begin{solution}{pb:b1:spherical-astronomy:1}
\textbf{1.} 余纬分别为 $\ang{90} - (\ang{-24.63}) = \ang{114.63}$ 与
$\ang{90} - \ang{51.48} = \ang{38.52}$；\omterm{def:b1:spherical-astronomy:geographic}{经度}差为 \ang{70.40}。

\textbf{2.} $\cos\Delta = \cos\ang{114.63}\cos\ang{38.52} +
\sin\ang{114.63}\sin\ang{38.52}\cos\ang{70.40} = -0.3260 + 0.1899 =
-0.1362$.

\textbf{3.} $\Delta = \ang{97.83}$；距离为 $6371\times1.7074 =
\qty{10878}{km}$，约 \qty{10900}{km}。

\textbf{4.}
\[
  \cos\mathrm{Az} = \frac{\cos\ang{38.52} - \cos\ang{114.63}\cos\ang{97.83}}
  {\sin\ang{114.63}\sin\ang{97.83}} = \frac{0.7823 - 0.0568}{0.9006} = 0.8056 ,
\]
$\mathrm{Az} = \ang{36.3}$，指向东北方向。

\textbf{5.} $\sin\mathrm{Az} = \sin\ang{70.40}\sin\ang{38.52}/\sin\ang{97.83}
= 0.5922$，由此得 \ang{36.3} 或 \ang{143.7}；
第 4 题中余弦为正，故取 \ang{36.3}。

\textbf{6.} 不能。等航向线与每条经线交成相同的角，
并向某一极点盘旋逼近；\omterm{def:b1:spherical-astronomy:spherical-triangle}{大圆}的航向则连续变化
（仅在起点处为 \ang{36.3}）。飞机若不随时修正航向，
将飞行更长的路程，并抵达别处。

\textbf{7.} 南\omterm{def:b1:spherical-astronomy:celestial-sphere}{天极}的高度
等于\omterm{def:b1:spherical-astronomy:geographic}{纬度}的绝对值：\ang{24.63}。

\textbf{8.} $h = \ang{90} - |\ang{-24.63} - (\ang{-16.72})| = \ang{82.09}$；
$\delta > \varphi$，故天狼星经过\omterm{def:b1:spherical-astronomy:celestial-sphere}{天顶}以北 \ang{7.91}。

\textbf{9.} 若 $\delta < -(\ang{90} - \ang{24.63}) =
\ang{-65.37}$，则恒星永不落下。老人星\omterm{def:b1:spherical-astronomy:equatorial}{赤纬}为 \ang{-52.70}，不是\omterm{def:b1:spherical-astronomy:circumpolar}{拱极星}，会落下。

\textbf{10.} $\ang{90} - |\ang{-24.63} + \ang{52.70}| = \ang{61.93}$，
位于\omterm{def:b1:spherical-astronomy:celestial-sphere}{天顶}以南。

\textbf{11.} $\cos H_0 = -\tan(\ang{-24.63})\tan(\ang{-16.72}) =
-(0.4584)(0.3004) = -0.1377$；$H_0 = \ang{97.92}$。天狼星
在地平线以上停留 $2\times97.92/15 = \qty{13.06}{h}$ 的\omterm{def:b1:spherical-astronomy:sidereal-time}{恒星时}，约 \qty{13}{h}。

\textbf{12.} $\cos\mathrm{A}_0 = \sin(\ang{-16.72})/\cos\ang{24.63} =
-0.3165$：天狼星在\omterm{def:b1:spherical-astronomy:horizontal}{方位角} \ang{108.5} 处升起，位于正东偏南。

\textbf{13.} 当 $H = 0$ 时，$\mathrm{LST} = \alpha = \qty{6}{h}~\qty{45}{min}~\qty{9}{s}$。

\textbf{14.} $\qty{86400}{s}\times365.2422/366.2422 = \qty{86164.09}{s} =
\qty{23}{h}~\qty{56}{min}~\qty{4}{s}$；经过时刻每夜提前
\qty{235.9}{s}，即 \qty{3}{min}~\qty{56}{s}。

\textbf{15.} 需要 $12 + n\times24/365.2422 = 6.7525 + 24$，故
$n = 18.7525/0.065710 = 285.4$ 天。自 3 月 20 日起 285 天为
12 月 30 日：天狼星在地方平子夜经过子午圈，发生于
2026 年 12 月 30 日至 31 日之夜。

\textbf{16.} 帕拉纳尔位于格林尼治以西 \ang{70.40}：地方平时
落后 $70.40/15 = \qty{4.693}{h}$，即 \qty{4}{h}~\qty{41.6}{min}。

\textbf{17.} \omterm{def:b1:spherical-astronomy:solar-time}{真太阳时} = \omterm{def:b1:spherical-astronomy:solar-time}{平太阳时} + \omterm{def:b1:spherical-astronomy:solar-time}{时差} =
$\qty{0}{h} - \qty{3}{min}$：日晷约指示 23:57。

\textbf{18.} $7\times\qty{3.93}{min} \approx \qty{27.5}{min}$，即提前。

\textbf{19.} 六月至日时为 $+\ang{23.44}$，
十二月至日时为 $-\ang{23.44}$。

\textbf{20.} 六月：$\ang{90} - |\ang{-24.63} - \ang{23.44}| = \ang{41.93}$。
十二月：$\ang{90} - |\ang{-24.63} + \ang{23.44}| = \ang{88.81}$；由于
$\delta_\odot = \ang{-23.44} > \varphi$，太阳经过
\omterm{def:b1:spherical-astronomy:celestial-sphere}{天顶}以北 \ang{1.19}。

\textbf{21.} $I = 1361\times(1/0.9837)^2\times\sin\ang{88.81} = 1361\times
1.0334\times0.9998 = \qty{1406}{W.m^{-2}}$。

\textbf{22.} 六月为南半球冬季。$\cos H_0 = -\tan(\ang{-24.63})
\tan\ang{23.44} = 0.4584\times0.4335 = 0.1988$：太阳在\omterm{def:b1:spherical-astronomy:equatorial}{时角}
$H_0 = \ang{78.53}$ 处落下。

\textbf{23.} $\cos\mathrm{A}_0 = \sin\ang{23.44}/\cos\ang{24.63} = 0.4376$：
$\mathrm{A}_0 = \ang{64.0}$，位于正东偏北。

\textbf{24.} 昼长 $2\times78.53/15 = \qty{10.47}{h}$；帕拉纳尔一年中最长的夜晚为
$24 - 10.47 = \qty{13.53}{h}$，即
\textbf{\qty[text-series-to-math=true]{13}{h}~\qty[text-series-to-math=true]{32}{min}}。
\end{solution}
